% ═══════════════════════════════════════════════════════════════
% SPANISCH MUSTERLÖSUNG ML-05a: LOS HOBBYS
% ═══════════════════════════════════════════════════════════════
% Fachfarbe: #c41e3a (Karminrot)
% Lösungsfarbe: ROT (deutlich sichtbar)
% Gesamtpunkte: 20P
% ═══════════════════════════════════════════════════════════════

\documentclass[11pt, a4paper]{article}

% ═══════════════════════════════════════════════════════════════
% SW-MODUS TOGGLE
% ═══════════════════════════════════════════════════════════════
\newif\ifswmodus
\swmodusfalse

% ═══════════════════════════════════════════════════════════════
% PAKETE
% ═══════════════════════════════════════════════════════════════

\usepackage[utf8]{inputenc}
\usepackage[T1]{fontenc}
\usepackage[ngerman]{babel}
\usepackage{helvet}
\renewcommand{\familydefault}{\sfdefault}

\usepackage[margin=2cm]{geometry}
\usepackage{enumitem}
\usepackage{tabularx}
\usepackage{booktabs}
\usepackage{multirow}

\usepackage{xcolor}
\usepackage{framed}
\usepackage{tcolorbox}
\tcbuselibrary{skins}

\usepackage{fancyhdr}
\usepackage{lastpage}
\usepackage{needspace}

% ═══════════════════════════════════════════════════════════════
% FARBEN
% ═══════════════════════════════════════════════════════════════

\definecolor{spanisch-color}{HTML}{c41e3a}
\definecolor{grau-color}{HTML}{888888}
\definecolor{hellgrau-color}{HTML}{f5f5f5}
\definecolor{orange-color}{HTML}{b35c00}
\definecolor{loesung-color}{HTML}{c62828}

\definecolor{sw-dark}{HTML}{000000}
\definecolor{sw-gray}{HTML}{666666}
\definecolor{sw-white}{HTML}{ffffff}

\ifswmodus
    \colorlet{spanisch}{sw-dark}
    \colorlet{grau}{sw-gray}
    \colorlet{hellgrau}{sw-white}
    \colorlet{orange}{sw-dark}
    \colorlet{loesung}{sw-dark}
\else
    \colorlet{spanisch}{spanisch-color}
    \colorlet{grau}{grau-color}
    \colorlet{hellgrau}{hellgrau-color}
    \colorlet{orange}{orange-color}
    \colorlet{loesung}{loesung-color}
\fi

\newcommand{\fachfarbe}{spanisch}

% ═══════════════════════════════════════════════════════════════
% VARIABLEN
% ═══════════════════════════════════════════════════════════════

\newcommand{\thema}{Los hobbys}
\newcommand{\sequenz}{05}
\newcommand{\block}{a}
\newcommand{\fach}{Spanisch}
\newcommand{\klasse}{7}
\newcommand{\maxpunkte}{20}
\newcommand{\datum}{\today}

% ═══════════════════════════════════════════════════════════════
% KOPF- UND FUSSZEILE
% ═══════════════════════════════════════════════════════════════

\pagestyle{fancy}
\fancyhf{}
\renewcommand{\headrulewidth}{0.4pt}
\renewcommand{\footrulewidth}{0.4pt}

\lhead{\fach{} -- Klasse \klasse}
\chead{\textcolor{loesung}{\textbf{ML-\sequenz\block{} (MUSTERLÖSUNG)}}}
\rhead{\datum}

\lfoot{}
\cfoot{}
\rfoot{Seite \thepage\ von \pageref{LastPage}}

% ═══════════════════════════════════════════════════════════════
% BEFEHLE
% ═══════════════════════════════════════════════════════════════

\newcommand{\punktebox}[1]{\hfill\fbox{\small \_/#1}}
\newcommand{\loesung}[1]{\textcolor{loesung}{\textbf{#1}}}
\newcommand{\luecke}[1]{\rule{#1}{0.4pt}}

\newtcolorbox{merkkasten}[1][]{
    colback=hellgrau,
    colframe=\fachfarbe,
    colbacktitle=white,
    coltitle=\fachfarbe,
    boxrule=1pt,
    arc=0pt,
    left=8pt, right=8pt, top=8pt, bottom=8pt,
    fonttitle=\bfseries,
    attach boxed title to top left={yshift=-2mm, xshift=5mm},
    boxed title style={boxrule=0pt, colback=white},
    title=#1
}

\newtcolorbox{vokabelkasten}{
    colback=white,
    colframe=\fachfarbe,
    boxrule=0.5pt,
    arc=0pt,
    left=5pt, right=5pt, top=5pt, bottom=5pt
}

\newtcolorbox{hinweis}{
    colback=white,
    colframe=orange,
    colbacktitle=white,
    coltitle=orange,
    boxrule=1pt,
    arc=0pt,
    left=5pt, right=5pt, top=5pt, bottom=5pt,
    fonttitle=\bfseries,
    attach boxed title to top left={yshift=-2mm, xshift=5mm},
    boxed title style={boxrule=0pt, colback=white},
    title=Hinweis
}

\newenvironment{aufgabe}[1]{%
    \needspace{5\baselineskip}%
    \nopagebreak[4]%
    \vspace{10pt}
    \noindent\textbf{\textcolor{\fachfarbe}{Aufgabe #1}}
    \nopagebreak[4]%
    \vspace{5pt}
    \nopagebreak[4]%
    \begin{list}{}{\leftmargin=0pt}
    \item[]
}{%
    \end{list}
}

\newcommand{\es}[1]{\textcolor{\fachfarbe}{\textbf{#1}}}

% ═══════════════════════════════════════════════════════════════
% DOKUMENT
% ═══════════════════════════════════════════════════════════════

\begin{document}

% ═══════════════════════════════════════════════════════════════
% KOPFBEREICH
% ═══════════════════════════════════════════════════════════════

\begin{center}
    {\Large\bfseries\textcolor{\fachfarbe}{Arbeitsblatt: \thema}}\\[0.3em]
    {\large\textcolor{loesung}{\textbf{--- MUSTERLÖSUNG ---}}}
\end{center}

\vspace{5pt}

\noindent
\begin{tabularx}{\textwidth}{|X|X|X|}
\hline
\textbf{Name:} \textcolor{loesung}{Musterschüler/in} &
\textbf{Klasse:} \klasse &
\textbf{Datum:} \datum \\
\hline
\end{tabularx}

\vspace{15pt}

% ═══════════════════════════════════════════════════════════════
% MERKKASTEN (mit Lösungen)
% ═══════════════════════════════════════════════════════════════

\begin{merkkasten}[Merkkasten: Hobbys / Los hobbys]
\textbf{Wichtige Vokabeln:}

\vspace{0.5em}

\begin{tabular}{ll}
\es{jugar al fútbol} & = \loesung{Fußball spielen} \\[0.8em]
\es{escuchar música} & = \loesung{Musik hören} \\[0.8em]
\es{ver la tele} & = \loesung{fernsehen} \\[0.8em]
\es{leer} & = \loesung{lesen} \\[0.8em]
\es{nadar} & = \loesung{schwimmen} \\[0.8em]
\es{montar en bici} & = \loesung{Fahrrad fahren} \\[0.8em]
\es{tocar la guitarra} & = \loesung{Gitarre spielen} \\[0.8em]
\es{chatear} & = \loesung{chatten} \\[0.8em]
\es{jugar a videojuegos} & = \loesung{Videospiele spielen} \\
\end{tabular}

\vspace{0.5em}

\textbf{Frage:} \es{¿Qué te gusta hacer?} = Was machst du gerne?

\textbf{Antwort:} \es{Me gusta} + Infinitiv (z.B. \textit{Me gusta nadar.})

\end{merkkasten}

\vspace{10pt}

% ═══════════════════════════════════════════════════════════════
% AUFGABE 1: Zuordnung (mit Lösungen)
% ═══════════════════════════════════════════════════════════════

\begin{aufgabe}{1}
Verbinde Spanisch und Deutsch. Ziehe Linien. \punktebox{4}
\end{aufgabe}

\begin{center}
\begin{tabular}{rcl}
\es{jugar al fútbol} & \loesung{---} & Gitarre spielen \loesung{(4)} \\[0.8em]
\es{escuchar música} & \loesung{---} & fernsehen \loesung{(3)} \\[0.8em]
\es{ver la tele} & \loesung{---} & Musik hören \loesung{(2)} \\[0.8em]
\es{tocar la guitarra} & \loesung{---} & Fußball spielen \loesung{(1)} \\
\end{tabular}
\end{center}

\vspace{0.5em}
\textit{\textcolor{grau}{Lösung: jugar al fútbol = Fußball spielen, escuchar música = Musik hören, ver la tele = fernsehen, tocar la guitarra = Gitarre spielen. Je 1P pro korrekter Zuordnung.}}

\vspace{15pt}

% ═══════════════════════════════════════════════════════════════
% AUFGABE 2: Lückentext (mit Lösungen)
% ═══════════════════════════════════════════════════════════════

\begin{aufgabe}{2}
Ergänze die Lücken mit den passenden Hobbys. \punktebox{5}
\end{aufgabe}

\begin{vokabelkasten}
\textit{Wortliste:} nadar | leer | chatear | montar en bici | jugar a videojuegos
\end{vokabelkasten}

\vspace{0.5em}

\noindent
a) Im Sommer gehe ich gerne ins Freibad. Auf Spanisch: Me gusta \loesung{nadar}.\\[0.8em]
b) Ich fahre jeden Tag mit dem Fahrrad zur Schule. Auf Spanisch: Me gusta \loesung{montar en bici}.\\[0.8em]
c) Abends schreibe ich mit meinen Freunden auf WhatsApp. Auf Spanisch: Me gusta \loesung{chatear}.\\[0.8em]
d) Am Wochenende spiele ich auf meiner Konsole. Auf Spanisch: Me gusta \loesung{jugar a videojuegos}.\\[0.8em]
e) Vor dem Einschlafen lese ich immer ein Buch. Auf Spanisch: Me gusta \loesung{leer}.\\

\vspace{0.5em}
\textit{\textcolor{grau}{Je 1P pro korrekter Einsetzung.}}

\vspace{15pt}

% ═══════════════════════════════════════════════════════════════
% AUFGABE 3: Dialog vervollständigen (mit Lösungen)
% ═══════════════════════════════════════════════════════════════

\begin{aufgabe}{3}
Vervollständige den Dialog mit deinen eigenen Antworten. \punktebox{5}
\end{aufgabe}

\noindent
\textbf{A:} \es{¡Hola! ¿Cómo te llamas?}

\textbf{B:} \loesung{Me llamo [Name]. / Soy [Name].}

\textbf{A:} \es{¿Qué te gusta hacer?}

\textbf{B:} \loesung{Me gusta [Hobby]. (z.B. Me gusta nadar.)}

\textbf{A:} \es{¡Qué bien! ¿Te gusta también nadar?}

\textbf{B:} \loesung{Sí, me gusta nadar. / No, no me gusta nadar.}

\textbf{A:} \es{¿Y qué más te gusta hacer?}

\textbf{B:} \loesung{Me gusta también [weiteres Hobby].}

\textbf{A:} \es{¡Muy interesante! ¡Hasta luego!}

\textbf{B:} \loesung{¡Hasta luego! / ¡Adiós!}

\vspace{0.5em}
\textit{\textcolor{grau}{Bewertung: 1P pro sinnvolle, sprachlich korrekte Antwort. Varianten möglich.}}

\vspace{15pt}

% ═══════════════════════════════════════════════════════════════
% AUFGABE 4: Freie Produktion (Beispiellösung)
% ═══════════════════════════════════════════════════════════════

\begin{aufgabe}{4}
Schreibe 3 Sätze über deine Hobbys. Verwende verschiedene Aktivitäten. \punktebox{6}
\end{aufgabe}

\textbf{Beispiellösung:}

\noindent
1. \loesung{Me gusta escuchar música.}

\vspace{0.5em}

\noindent
2. \loesung{Me gusta jugar al fútbol.}

\vspace{0.5em}

\noindent
3. \loesung{Me gusta ver la tele.}

\vspace{0.5em}
\textit{\textcolor{grau}{Bewertung pro Satz: Struktur ,,Me gusta + Infinitiv'' korrekt (1P), Hobby korrekt geschrieben (1P). = 2P pro Satz, max. 6P. Verschiedene Hobbys verwenden!}}

% ═══════════════════════════════════════════════════════════════
% PUNKTEÜBERSICHT
% ═══════════════════════════════════════════════════════════════

\vspace{20pt}
\hrule
\vspace{10pt}

\begin{flushright}
\textbf{Punkte:} \loesung{\maxpunkte} / \maxpunkte
\end{flushright}

\end{document}
